% !TEX program = xelatex
\documentclass[11pt]{ctexart}
%=====================================================================
% theme-violet.tex —— 讲义精讲 · 紫花主题
% 用法：\documentclass[11pt]{ctexart}
%       \input{assets/theme-violet}
%       \begin{document} ... \end{document}
% 只用 ctexart 自带能力 + xcolor/enumitem/fancyhdr/pifont/amsmath，
% 不依赖 tikz/tcolorbox，保证 MiKTeX 可直接编译。
%=====================================================================

\usepackage[a4paper,left=2.0cm,right=2.0cm,top=2.1cm,bottom=2.3cm,%
            headsep=12pt,footskip=22pt]{geometry}
\usepackage{amsmath,amssymb}
\usepackage{xcolor}
\usepackage{array}
\usepackage{enumitem}
\usepackage{pifont}
\usepackage{fancyhdr}
\usepackage[hidelinks]{hyperref}

%---------------------------------------------------------------- 配色
\definecolor{violet}{RGB}{107,78,143}      % 主色：紫
\definecolor{violetdeep}{RGB}{63,44,92}    % 深紫：标题/强调
\definecolor{petal}{RGB}{183,159,208}      % 花瓣：装饰线、花饰
\definecolor{bandviolet}{RGB}{247,243,251} % 浅紫底：档案/提示框
\definecolor{bandcream}{RGB}{252,249,242}  % 米白底：易错/注意框
\definecolor{rulepurple}{RGB}{217,207,230} % 分隔线
\definecolor{inkgray}{RGB}{92,96,110}      % 次要文字

%---------------------------------------------------------------- 字体
% 有 Noto CJK 用 Noto，没有就退回 ctex 默认（Windows 下为宋体/黑体）
\IfFontExistsTF{Noto Serif CJK SC}{\setCJKmainfont{Noto Serif CJK SC}}{}
\IfFontExistsTF{Noto Sans CJK SC}{\setCJKsansfont{Noto Sans CJK SC}}{}

%---------------------------------------------------------------- 版式
\setlist{nosep,leftmargin=1.9em}
\linespread{1.28}
\setlength{\parskip}{3.5pt}
\setlength{\parindent}{0pt}
\setlength{\tabcolsep}{5pt}
\renewcommand{\arraystretch}{1.28}

%---------------------------------------------------------------- 花饰
% \ding{95} = 五瓣小花，\ding{96} = 星形花；都来自 pifont，必定存在
\newcommand{\orn}{\textcolor{petal}{\ding{95}}}          % 行内小标
\newcommand{\ornstar}{\textcolor{petal}{\ding{96}}}      % 星形花
\newcommand{\flowerrule}{%                                % 居中花饰分隔线
  \par\vspace{7pt}%
  \noindent\makebox[\linewidth]{\color{petal}%
    \rule{0.20\linewidth}{0.6pt}\hspace{8pt}{\small\ding{95}}\hspace{4pt}%
    \raisebox{1.1pt}{\tiny\ding{95}}\hspace{4pt}{\small\ding{95}}\hspace{8pt}%
    \rule{0.20\linewidth}{0.6pt}}%
  \par\vspace{7pt}}

%---------------------------------------------------------------- 标题与节头
\newcommand{\doctitle}[3]{%   #1 主标题  #2 副标题  #3 元信息（可用 \\ 分行）
  \begin{center}
    {\color{petal}\small \orn\enspace #2\enspace \ornstar}\\[7pt]
    {\LARGE\bfseries\color{violetdeep} #1}\\[10pt]
    {\color{petal}\rule{0.55\linewidth}{1pt}}\\[4pt]
    {\small\color{inkgray} #3}
  \end{center}\vspace{2pt}}

\newcommand{\hdr}[2]{%   #1 章节号（如 §0）  #2 标题
  \par\vspace{12pt}%
  \noindent\colorbox{violetdeep}{\color{white}\bfseries\small\;#1\;}\hspace{7pt}%
  {\color{violetdeep}\large\bfseries #2}\par\nobreak
  \vspace{3pt}{\color{petal}\hrule height 1.2pt}\vspace{6pt}}

\newcommand{\subhdr}[1]{%
  \par\vspace{6pt}%
  \noindent{\color{violet}\bfseries #1}\par\vspace{1pt}}

%---------------------------------------------------------------- 提示框
\newcommand{\callout}[1]{% 浅紫底：讲义档案、读法说明
  \par\medskip\noindent
  \fcolorbox{rulepurple}{bandviolet}{\parbox{\dimexpr\linewidth-2\fboxsep-2\fboxrule}{\vspace{2pt}#1\vspace{2pt}}}%
  \par\medskip}

\newcommand{\warnbox}[1]{% 米白底：易错、边界条件
  \par\medskip\noindent
  \fcolorbox{rulepurple}{bandcream}{\parbox{\dimexpr\linewidth-2\fboxsep-2\fboxrule}{\vspace{2pt}#1\vspace{2pt}}}%
  \par\medskip}

\newcommand{\gapitem}[2]{% #1 缺口类型（紫） #2 内容
  \par\vspace{2pt}\noindent
  \hangindent=1.4em\hangafter=1
  {\color{petal}\small\ding{96}}\enspace{\bfseries\color{violetdeep}#1}\enspace #2\par}

% 前提旁注：只放在「需要理解的那一步」后面，给陈述 + 本题代入，不重证。
\newcommand{\aside}[1]{%
  \par\vspace{1pt}\noindent
  \hangindent=1.6em\hangafter=1
  {\small\color{inkgray}%
    \hspace*{0.6em}{\color{petal}$\triangleright$}\hspace{2pt}#1}%
  \par\vspace{2pt}}

%---------------------------------------------------------------- 表格
% 本机 MiKTeX 的 colortbl 与 array 版本不匹配（\rowcolor 会报错），
% 故表格统一用「上粗紫线 + 表头 + 中细线」的三线式，不依赖 colortbl。
\newcommand{\tabtop}{\par\vspace{3pt}{\color{violet}\hrule height 1.2pt}\vspace{1pt}}
\newcommand{\tabmid}{\par\vspace{1pt}{\color{petal}\hrule height 0.5pt}\vspace{1pt}}
\newcommand{\tabbot}{\par\vspace{3pt}{\color{rulepurple}\hrule height 0.5pt}\vspace{8pt}}
\newcommand{\thead}[1]{{\bfseries\color{violetdeep}#1}}

%---------------------------------------------------------------- 页眉页脚
\pagestyle{fancy}
\fancyhf{}
\renewcommand{\headrulewidth}{0.4pt}
\renewcommand{\headrule}{\hbox to\headwidth{\color{rulepurple}\leaders\hrule height \headrulewidth\hfill}}
\fancyhead[L]{\small\color{inkgray}\songti \doctag}
\fancyhead[R]{\small\color{inkgray}\songti \coursetag}
\fancyfoot[C]{\small\color{inkgray}\orn\enspace{\color{violet}\thepage}\enspace\orn}

% 每篇讲义在导言区设这两个值
\newcommand{\doctag}{}
\newcommand{\coursetag}{}
\newcommand{\settags}[2]{\renewcommand{\doctag}{#1}\renewcommand{\coursetag}{#2}}

%---------------------------------------------------------------- 数学记号
\newcommand{\vect}[1]{\boldsymbol{#1}}
\newcommand{\dd}{\,\mathrm{d}}
\newcommand{\ddx}{\frac{\mathrm{d}}{\mathrm{d}x}}
\addtolength{\jot}{2pt}



\setlength{\headheight}{22pt}
\usepackage{tikz}
\usetikzlibrary{positioning,arrows.meta,calc}
\usepackage[american]{circuitikz}
\settags{实验 4 · 直流稳压电源}{电工电子技术}
\let\orighdr\hdr
\renewcommand{\hdr}[2]{\needspace{12\baselineskip}\orighdr{#1}{#2}}
\let\origsubhdr\subhdr
\renewcommand{\subhdr}[1]{\needspace{3\baselineskip}\origsubhdr{#1}}

\begin{document}

\doctitle{4.\ 直流稳压电源实验}
{从 220\,V 交流到 12\,V 直流：变压、整流、滤波、稳压四步怎么走、怎么算、怎么判}
{依据：实验手册 pp.10--13（Experiment 4, DC Stabilized Power Supply）；课程课件 C3-2（整流，pp.13--25）与 C3-3（滤波与稳压，pp.8--24）；教材《电工学实验教程》第 2 版有关测量与误差的约定\ \textperiodcentered\ 2026-10-27}

\callout{\textbf{这份讲义怎么用。} §2 讲清楚“为什么会有 $0.9U_2$ 和 $1.2U_2$”，§3 把实验台的电路拆成三条支路，§5 是按顺序可照做的流程，§6--§7 是数据表与逐项分析，§8 是手册 p13 五道 Discussion 的参考答法。所有以 $U_2=14.0\,\mathrm{V}$ 开头的数字都只是\textbf{算例}：进实验室后先用万用表测出你自己变压器副边的 $U_2$，再把算例里的 14.0\,V 换掉。}

%=============================================================
\hdr{§0}{手册坐标、记号与一处必须更正的开关状态}

\subhdr{0.1 手册这四页各是什么}

\begin{center}
\small
\begin{tabular}{@{}p{1.7cm}p{5.2cm}p{8.4cm}@{}}
\hline
\thead{页} & \thead{手册内容} & \thead{本讲义对应} \\
\hline
p10 & Preliminary Work：4 个预习问题；预习 DS2202A 示波器 & §4 逐题回答；§5.8 示波器速查 \\
p11 & Experimental Procedure：实验电路与桥式整流、电容滤波的测量步骤 & §2 原理；§3 电路；§5.3--§5.5 流程 \\
p12 & Table 1、Table 2；稳压电路的测量步骤 & §6 数据表；§5.6 稳压测量 \\
p13 & Discussion of Results：5 个讨论题 & §7 分析；§8 逐题参考答法 \\
\hline
\end{tabular}
\end{center}

\subhdr{0.2 记号与单位：先统一，后文不再回头}

\begin{center}
\small
\begin{tabular}{@{}p{2.2cm}p{8.8cm}p{4.0cm}@{}}
\hline
\thead{符号} & \thead{含义} & \thead{单位} \\
\hline
$U_2,\ u_2$ & 变压器副边电压的有效值 / 瞬时值 & V \\
$U_{o1}$ & 整流或滤波后的输出平均值（直流分量） & V \\
$U_{o2}$ & 稳压管稳压后的输出平均值 & V \\
$U_Z$ & 稳压管的稳压值（本实验 $12\,\mathrm{V}$） & V \\
$I_o,\ I_L$ & 负载电流 & mA \\
$I_D$ & 桥式整流中每只二极管的平均电流 & mA \\
$I_R,\ I_Z$ & 限流电阻 $R_2$ 的电流 / 稳压管电流 & mA \\
$R_L$ & 负载电阻（表 1 用 $R_1$，表 2 用 $R_w+R_3$） & $\mathrm{k}\Omega$ \\
$f,\ T$ & 电网频率 / 周期 & $50\,\mathrm{Hz}$ / $20\,\mathrm{ms}$ \\
\hline
\end{tabular}
\end{center}

\aside{万用表的交流挡读的是\textbf{有效值}，直流挡读的是\textbf{平均值}，示波器读的是\textbf{瞬时值}（峰峰值、峰值、周期）。同一个量用不同仪表读，数值不完全相同是正常的；报告里要写清“用什么挡、读的是什么值”。}

\subhdr{0.3 手册 p11 第三步的开关状态：按电路改正}

手册 p11 第三步原文是：“Disconnect K1 and K2, connect K3. Measure $U_{o1}$ across the $5.1\,\mathrm{k}\Omega$ resistor with a multimeter.” 这句话与电路图对不上：$R_1=5.1\,\mathrm{k}\Omega$ 在 $K_2$ 支路里，$K_2$ 断开后 $R_1$ 已经脱离电路；$K_3$ 接通的是 $R_2$、稳压管和 $R_w+R_3$ 那一条支路，量不到“$5.1\,\mathrm{k}\Omega$ 上的 $U_{o1}$”。Table 1 第三行明确写的是“$U_{o1}$ Output voltage of capacitor filter circuit”，所以这一步应当是让 $C_1$ 与 $R_1$ 同时接入，即 $K_1$、$K_2$ 闭合，$K_3$ 断开。按下面的真值表接线，才不会出现“表上读数为零、波形也不对”的假故障。若老师现场对开关编号另有说明，以老师为准；但从电路和 Table 1 的定义看，测滤波必须让 $C_1$ 与 $R_1$ 同时接入。

\begin{center}
\small
\begin{tabular}{@{}p{3.6cm}p{1.1cm}p{1.1cm}p{1.1cm}p{5.4cm}@{}}
\hline
\thead{测量内容} & \thead{$K_1$} & \thead{$K_2$} & \thead{$K_3$} & \thead{输出端接什么} \\
\hline
变压器副边 $U_2$ & 断 & 断 & 断 & 副边两端，交流电压 \\
桥式整流 $U_{o1}$ & 断 & 合 & 断 & $R_1=5.1\,\mathrm{k}\Omega$ 两端 \\
电容滤波 $U_{o1}$ & 合 & 合 & 断 & $R_1$ 两端，$C_1$ 并联在旁 \\
稳压输出 $U_{o2}$ & 合 & 断 & 合 & $R_w+R_3$ 两端 \\
\hline
\end{tabular}
\end{center}

\warnbox{\textbf{一条硬规则。} 手册 p11 的 Caution 写明：接线、换元件、拆线之前，必须先关变压器模块的电源开关。改接 $K_1$--$K_3$ 时不要带电操作；$C_1$ 断电后仍可能存电，改接前用电阻放电。}

%=============================================================
\hdr{§1}{实验目的与安全}

\subhdr{1.1 这次实验要完成的五件事}

\begin{enumerate}
  \item 用万用表交流挡测变压器副边电压 $U_2$，用示波器观察并画出副边正弦波形。
  \item 只接 $R_1$（$K_2$ 合，$K_1$、$K_3$ 断），测桥式整流输出 $U_{o1}$，观察全波整流波形，验证 $U_{o1}\approx0.9U_2$。
  \item 接入 $C_1$（$K_1$、$K_2$ 合，$K_3$ 断），测滤波输出 $U_{o1}$，观察带纹波的直流波形，验证 $U_{o1}\approx1.2U_2$（带负载）。
  \item 接入稳压支路（$K_1$、$K_3$ 合，$K_2$ 断），分别在 $R_w+R_3=3\,\mathrm{k}\Omega$ 与 $103\,\mathrm{k}\Omega$ 两种负载下测 $U_{o2}$，判断稳压管是否工作在击穿区。
  \item 用测得的数据完成 Table 1、Table 2 与 p13 的五道 Discussion；所有计算题写出公式、代入数值、给出单位和结论。
\end{enumerate}

\subhdr{1.2 安全：220\,V 在一次侧，低压不等于无风险}

\begin{itemize}
  \item 变压器一次侧接 220\,V 市电。\textbf{不得触碰一次侧接线柱、插头和裸露导线}；通电前请老师检查接线。
  \item 改接任何导线前先关变压器模块电源；单手操作，站在绝缘垫上，不要两个人同时碰电路的不同部位。
  \item $C_1$ 是电解电容，断电后仍可能存电。改接或拆线前用 $1\,\mathrm{k}\Omega$ 左右的电阻跨接其两端放电，不要用螺丝刀直接短路。
  \item 稳压管、整流桥都有极性；$C_1$ 的反接会发热、鼓包甚至爆裂。接线前核对极性，通电后先看有无异味、冒烟、异常发热。
  \item 示波器探头的地线夹必须接电路的公共地（负极），不能夹在带电节点上；否则会通过示波器地线造成短路。
  \item 发现保险丝熔断、二极管发烫、输出电压为零等异常，立即关电源并报告老师，不要反复通电“再试一次”。
\end{itemize}

%=============================================================
\hdr{§2}{原理：从 220\,V 交流到 12\,V 直流}

\subhdr{2.1 四步链：变压—整流—滤波—稳压}

\begin{figure}[htbp]
\centering
\begin{tikzpicture}[node distance=0.75cm,
  box/.style={draw=violet, rounded corners=2pt, minimum width=2.0cm,
  minimum height=1.0cm, align=center, fill=bandviolet, font=\small}]
\node[box] (T) {变压器\\$T_1$};
\node[box, right=of T] (B) {桥式整流\\$D_1$};
\node[box, right=of B] (C) {电容滤波\\$C_1$};
\node[box, right=of C] (R) {稳压管\\$R_2+D_Z$};
\node[box, right=of R] (L) {负载\\$R_w+R_3$};
\draw[-{Latex}] (T) -- node[above,font=\scriptsize]{$u_2$} (B);
\draw[-{Latex}] (B) -- node[above,font=\scriptsize]{$U_{o1}$} (C);
\draw[-{Latex}] (C) -- node[above,font=\scriptsize]{$U_{o1}$} (R);
\draw[-{Latex}] (R) -- node[above,font=\scriptsize]{$U_{o2}$} (L);
\node[below=0.15cm of T, font=\scriptsize] {220\,V, 50\,Hz};
\end{tikzpicture}
\caption{直流稳压电源的四步链。变压器只改变交流电压的幅值；整流把负半周翻正；滤波把脉动削平；稳压管把输出钉在 $U_Z$ 附近。}
\end{figure}

这四步的分工可以用一句话记住：\textbf{变压器变幅值，整流变方向，滤波变平滑，稳压变准确}。下面逐级说明每一步的输出是什么、为什么是那个数。

\subhdr{2.2 变压器：只变幅值，不改变频率}

变压器副边电压仍是有正有负的正弦交流：
\[
u_2(t)=\sqrt2\,U_2\sin(2\pi ft),\qquad f=50\,\mathrm{Hz},\quad T=20\,\mathrm{ms}.
\]
若一次侧为 $U_1$、匝数为 $N_1$，副边匝数为 $N_2$，则理想情况下
\[
\frac{U_1}{U_2}=\frac{N_1}{N_2}.
\]
本实验的变压器模块把 220\,V 降到低压（具体数值以模块铭牌和实测为准），\textbf{频率仍是 50\,Hz}，波形仍是正弦。注意：变压器不整流，副边电压的平均值为零；用万用表直流挡去量副边，读数应接近零，这是正常的。

\subhdr{2.3 桥式整流：把负半周翻上来}

桥式整流由四只二极管 $D_1$--$D_4$ 组成。副边正半周时，一对二极管导通；负半周时，另一对导通。无论哪个半周，流过负载 $R_1$ 的电流方向都不变，所以输出是\textbf{全波脉动直流}。

\begin{center}
\small
\begin{tabular}{@{}p{2.0cm}p{2.8cm}p{9.0cm}p{1.8cm}@{}}
\hline
\thead{半周} & \thead{导通的二极管} & \thead{电流路径（从副边看）} & \thead{输出} \\
\hline
$u_2>0$ & $D_1,\ D_4$ & 副边上端 $\to D_1\to R_1\to D_4\to$ 副边下端 & 正半周 \\
$u_2<0$ & $D_2,\ D_3$ & 副边下端 $\to D_2\to R_1\to D_3\to$ 副边上端 & 翻正后的半周 \\
\hline
\end{tabular}
\end{center}

把 $|\sin\omega t|$ 在一个周期内取平均，得到全波整流的输出平均值：
\[
U_{o1}=\frac{1}{\pi}\int_0^{\pi}\sqrt2\,U_2\sin\theta\,d\theta
=\frac{2\sqrt2}{\pi}U_2\approx0.9\,U_2.
\]
这个 $0.9$ 不是经验系数，而是“全波整流正弦波的平均值”算出来的；若考虑两只二极管的正向压降，实际值会略低，约 $0.9U_2-2U_D$。

整流输出的特点是：
\begin{itemize}
  \item 方向不变，但大小仍在脉动，所以叫“脉动直流”，不是稳定直流；
  \item 脉动频率是电网频率的 2 倍，即 $100\,\mathrm{Hz}$，周期 $10\,\mathrm{ms}$；
  \item 每只二极管只承担一半的负载电流：$I_D=I_o/2$，其中 $I_o=U_{o1}/R_L$；
  \item 每只二极管承受的最高反向电压为 $U_{Rm}=\sqrt2\,U_2$。
\end{itemize}

\subhdr{2.4 电容滤波：把脉动削平}

把 $C_1$ 并联在负载两端后，整流输出给电容充电；当整流电压低于电容电压时，二极管截止，电容通过负载放电。电容反复“充电—放电”，把电压抬到接近峰值，并把谷值填高，输出就变成带少量纹波的直流。

带负载时，工程上常用
\[
U_{o1}\approx1.2\,U_2
\]
作为估算值；空载时电容充到峰值，$U_{o1}\approx\sqrt2\,U_2\approx1.4\,U_2$。所以 $1.2$ 与 $1.4$ 并不矛盾：$1.2$ 是“带负载、有放电”的结果，$1.4$ 是“空载、几乎不放电”的结果。实际值取决于时间常数
\[
\tau=R_LC_1.
\]
$\tau$ 越大（$C_1$ 越大或负载越轻），放电越慢，输出越接近峰值，纹波越小。

纹波的大小可以用一个常用估算式判断：
\[
V_r\approx\frac{I_o}{2fC_1},\qquad I_o=\frac{U_{o1}}{R_L}.
\]
本实验 $C_1=47\,\mu\mathrm{F}$、$R_1=5.1\,\mathrm{k}\Omega$、$f=50\,\mathrm{Hz}$，若 $U_{o1}\approx16.8\,\mathrm{V}$，则 $I_o\approx3.29\,\mathrm{mA}$，
\[
V_r\approx\frac{3.29\times10^{-3}}{2\times50\times47\times10^{-6}}\,\mathrm{V}\approx0.70\,\mathrm{V}.
\]
与 $16.8\,\mathrm{V}$ 的直流分量相比，纹波约 $4\%$，所以示波器上看到的是一条“略有锯齿”的直流线，而不是完全平直。

\begin{figure}[htbp]
\centering
\begin{tikzpicture}[scale=0.95,>=Latex]
% (a)
\begin{scope}[xshift=0cm]
\draw[->] (0,0) -- (4.3,0) node[right] {$\omega t$};
\draw[->] (0,-1.7) -- (0,1.9) node[above] {$u_2$};
\draw[thick,blue,domain=0:720,samples=200] plot ({\x/180*0.9},{1.414*sin(\x)});
\draw[dashed] (0,1.414) -- (3.9,1.414);
\node[above] at (2.0,1.45) {$\sqrt2 U_2$};
\foreach \x in {0.9,1.8,2.7,3.6} \draw (\x,0.05) -- (\x,-0.05);
\node[below] at (0.9,-0.05) {$\pi$};
\node[below] at (1.8,-0.05) {$2\pi$};
\node[below] at (2.7,-0.05) {$3\pi$};
\node[below] at (3.6,-0.05) {$4\pi$};
\node[below left] at (0,0) {0};
\node at (2.0,-2.0) {(a) 副边正弦};
\end{scope}
% (b)
\begin{scope}[xshift=5.6cm]
\draw[->] (0,0) -- (4.3,0) node[right] {$\omega t$};
\draw[->] (0,-0.2) -- (0,1.9) node[above] {$u_{o1}$};
\draw[thick,red,domain=0:720,samples=200] plot ({\x/180*0.9},{1.414*abs(sin(\x))});
\draw[dashed] (0,0.9) -- (3.9,0.9);
\node[font=\scriptsize] at (2.0,1.65) {虚线：平均值 $0.9U_2$};
\foreach \x in {0.9,1.8,2.7,3.6} \draw (\x,0.05) -- (\x,-0.05);
\node[below] at (0.9,-0.05) {$\pi$};
\node[below] at (1.8,-0.05) {$2\pi$};
\node[below] at (2.7,-0.05) {$3\pi$};
\node[below] at (3.6,-0.05) {$4\pi$};
\node[below left] at (0,0) {0};
\node at (2.0,-2.0) {(b) 桥式整流};
\end{scope}
% (c)
\begin{scope}[xshift=11.2cm]
\draw[->] (0,0) -- (4.3,0) node[right] {$\omega t$};
\draw[->] (0,-0.2) -- (0,1.9) node[above] {$u_{o1}$};
\draw[thick,red] (0,1.05) -- (0.10,1.38) -- (0.35,1.32) -- (0.70,1.22) -- (0.95,1.10)
 -- (1.05,1.38) -- (1.30,1.32) -- (1.65,1.22) -- (1.90,1.10)
 -- (2.00,1.38) -- (2.25,1.32) -- (2.60,1.22) -- (2.85,1.10)
 -- (2.95,1.38) -- (3.20,1.32) -- (3.55,1.22) -- (3.80,1.10)
 -- (3.90,1.38);
\draw[dashed] (0,1.2) -- (3.9,1.2);
\node[font=\scriptsize] at (2.0,1.65) {虚线：平均值 $1.2U_2$};
\draw[<->] (2.1,1.10) -- (2.1,1.38);
\node[fill=white,inner sep=1pt,font=\scriptsize] at (2.3,1.24) {$V_r$};
\node at (2.0,-2.0) {(c) 电容滤波};
\end{scope}
\end{tikzpicture}
\caption{三个关键波形（纵轴以 $U_2$ 为基准，纹波已适当放大）：(a) 副边正弦，峰值 $\sqrt2U_2$、周期 $20\,\mathrm{ms}$；(b) 全波整流，平均值 $0.9U_2$、脉动周期 $10\,\mathrm{ms}$；(c) 电容滤波，平均值约 $1.2U_2$，峰谷差为纹波 $V_r$。}
\end{figure}

\subhdr{2.5 稳压管：把输出钉在 $U_Z$ 附近}

稳压管 $D_Z$ 反向接在负载两端，$R_2$ 串在整流滤波输出与稳压管之间，起限流和分压作用。当 $D_Z$ 进入反向击穿区时，它两端的电压近似为 $U_Z$（本实验 $12\,\mathrm{V}$），只要流过它的电流仍在 $I_{Z\min}$ 与 $I_{Z\max}$ 之间，输出就基本不变。

对稳压支路写节点电流方程：
\[
I_R=\frac{U_{o1}-U_Z}{R_2},\qquad
I_L=\frac{U_Z}{R_L},\qquad
I_Z=I_R-I_L.
\]
判断稳压管是否在工作，不能只看“输出电压是不是 12\,V”，而要按下面三步算：
\begin{enumerate}
  \item 先假设稳压管击穿、$U_{o2}=U_Z$；
  \item 算出 $I_R$ 与 $I_L$，得到 $I_Z=I_R-I_L$；
  \item 若 $I_{Z\min}\le I_Z\le I_{Z\max}$，假设成立，稳压管在工作；若 $I_Z<0$（或小于 $I_{Z\min}$），说明负载“抢走”了全部电流，稳压管退出击穿，此时要按分压重新算：
  \[
  U_{o2}=\frac{R_L}{R_2+R_L}\,U_{o1}.
  \]
\end{enumerate}
若 $I_Z>I_{Z\max}$，稳压管会过热甚至损坏，必须增大 $R_2$、降低 $U_{o1}$，或加一个泄放电阻保证最小负载电流；不能靠“减轻负载”来解决，因为负载越轻，$I_L$ 越小，$I_Z$ 反而越大。$R_2$ 因此被称为\textbf{限流电阻}：它既要在最重负载时给稳压管留一点电流，又要在最轻负载时限住稳压管的电流。

\subhdr{2.6 整条链的预期值：以 $U_2=14.0\,\mathrm{V}$ 为例}

\begin{center}
\small
\begin{tabular}{@{}p{4.8cm}p{4.2cm}p{5.6cm}@{}}
\hline
\thead{测量点} & \thead{估算值} & \thead{依据} \\
\hline
副边 $U_2$ & $14.0\,\mathrm{V}$（实测为准） & 变压器铭牌/实测 \\
桥式整流 $U_{o1}$ & $0.9\times14.0=12.6\,\mathrm{V}$ & 全波整流平均值 \\
电容滤波 $U_{o1}$（带 $5.1\,\mathrm{k}\Omega$） & $1.2\times14.0=16.8\,\mathrm{V}$ & 带负载经验式 \\
空载滤波 $U_{o1}$ & $\sqrt2\times14.0\approx19.8\,\mathrm{V}$ & 电容充到峰值 \\
稳压 $U_{o2}$（$R_L=103\,\mathrm{k}\Omega$） & $12\,\mathrm{V}$ & 稳压管工作 \\
稳压 $U_{o2}$（$R_L=3\,\mathrm{k}\Omega$） & 约 $0.6\times16.8=10.1\,\mathrm{V}$ & 稳压管退出，按分压 \\
\hline
\end{tabular}
\end{center}

这张表只是“进实验室前心里有数”的预算，不是让你照抄。若你的 $U_2$ 是 12\,V 或 15\,V，把 14.0 换掉重算即可；定性结论（整流约 0.9 倍、滤波约 1.2 倍、重载时稳压管可能退出）不变。

%=============================================================
\hdr{§3}{实验电路与三条支路}

\subhdr{3.1 把手册 p11 的图拆成三条并联支路}

手册 p11 的电路可以这样读：变压器 $T_1$ 的副边先送进桥堆 $D_1$；桥堆的输出上端是正母线，下端是负母线。正母线与负母线之间挂了三条可以独立投切的支路：
\begin{itemize}
  \item $K_1$ 支路：$C_1=47\,\mu\mathrm{F}$ 滤波电容；
  \item $K_2$ 支路：$R_1=5.1\,\mathrm{k}\Omega$ 负载电阻；
  \item $K_3$ 支路：$R_2=2\,\mathrm{k}\Omega$ 限流电阻，后面接稳压管 $D_Z$ 与负载 $R_w+R_3$。
\end{itemize}
$U_{o1}$ 是正母线与负母线之间的电压（也就是 $C_1$ 或 $R_1$ 两端）；$U_{o2}$ 是稳压管两端（也就是 $R_w+R_3$ 两端）的电压。

\begin{figure}[htbp]
\centering
\begin{circuitikz}[scale=0.9,transform shape]
\node[draw,minimum width=1.9cm,minimum height=1.0cm,align=center] (T) at (0,2) {T1\\[1pt]{\small 220\,V/低压}};
\node[draw,minimum width=1.9cm,minimum height=1.0cm,align=center] (B) at (3.6,2) {D1\\[1pt]{\small bridge}};
\draw (T) -- (B);
\draw (B.north) -- (11,2.5);
\draw (B.south) -- (3.6,0) -- (16,0);
\draw (7,2.5) to[nos] (7,1.6) to[C] (7,0);
\node[right] at (7.15,2.05) {$K_1$};
\node[right] at (7.15,0.75) {$C_1$};
\draw (9,2.5) to[nos] (9,1.6) to[R] (9,0);
\node[right] at (9.15,2.05) {$K_2$};
\node[right] at (9.15,0.75) {$R_1$};
\draw (11,2.5) to[nos] (11,1.6) to[R] (11,0.8) -- (13,0.8);
\node[right] at (11.15,2.05) {$K_3$};
\node[right] at (11.15,1.2) {$R_2$};
\draw (13,0.8) to[zD] (13,0);
\node[left] at (12.7,0.4) {$D_Z$};
\draw (13,0.8) to[vR, l=$R_w$] (14.2,0.8) to[R, l=$R_3$] (15.4,0.8) -- (16.1,0.8);
\draw (15.4,0) -- (16.1,0);
\draw[<->] (6.2,2.4) -- (6.2,0.1) node[midway,left] {$U_{o1}$};
\draw[<->] (16.1,0.7) -- (16.1,0.1) node[midway,right] {$U_{o2}$};
\end{circuitikz}
\caption{测量电路示意（简化）：$T_1$ 与 $D_1$ 用方框表示，重点是 $K_1$--$K_3$ 三条支路和 $U_{o1}$、$U_{o2}$ 的测量位置。完整接线见手册 p11。}
\end{figure}

\subhdr{3.2 元件与作用}

\begin{center}
\small
\begin{tabular}{@{}p{1.8cm}p{3.2cm}p{9.2cm}@{}}
\hline
\thead{元件} & \thead{参数} & \thead{作用} \\
\hline
$T_1$ & 220\,V/低压 & 把市电降到安全低压，只变幅值、不变频率 \\
$D_1$ & 桥堆（4 只二极管） & 全波整流，把负半周翻正 \\
$C_1$ & $47\,\mu\mathrm{F}$ & 滤波，减小纹波；同时是“储能”元件，断电后要放电 \\
$R_1$ & $5.1\,\mathrm{k}\Omega$ & 整流/滤波实验的负载，保证有电流可测 \\
$R_2$ & $2\,\mathrm{k}\Omega$ & 稳压支路的限流电阻，限制 $I_Z$ \\
$D_Z$ & $12\,\mathrm{V}$ 稳压管 & 反向击穿时把 $U_{o2}$ 钉在 $12\,\mathrm{V}$ 附近 \\
$R_w$ & $0$--$100\,\mathrm{k}\Omega$ 可调 & 与 $R_3$ 串联，改变稳压支路的负载 \\
$R_3$ & $3\,\mathrm{k}\Omega$ & 与 $R_w$ 一起构成稳压支路的负载 \\
$K_1$--$K_3$ & 开关 & 分别投切 $C_1$、$R_1$、稳压支路 \\
\hline
\end{tabular}
\end{center}

\subhdr{3.3 三条支路的投切与测量点}

把 §0.3 的真值表再具体化一次：测 $U_2$ 时三条支路全断；测整流时只合 $K_2$；测滤波时合 $K_1$、$K_2$；测稳压时合 $K_1$、$K_3$，并把 $K_2$ 断开。这样 $U_{o1}$ 始终是“正母线与负母线”之间的电压，$U_{o2}$ 始终是“稳压管两端”的电压，不会因为支路投错而读到 0 或悬浮值。

\warnbox{\textbf{测量前先看三件事。} （1）万用表挡位：$U_2$ 用交流电压挡，$U_{o1}$、$U_{o2}$ 用直流电压挡；（2）示波器耦合：$U_2$ 用 AC 耦合，整流/滤波/稳压输出用 DC 耦合；（3）共地：示波器探头地线夹接负母线，不要接在正母线或开关两端。}

%=============================================================
\hdr{§4}{预习问题（手册 p10）}

\subhdr{4.1 变压器副边电压是直流还是交流？画出输出波形}

副边是\textbf{交流}。变压器靠电磁感应工作，副边电压与一次侧同频率、同波形，只是幅值变了；它仍有正负半周，平均值在一个周期内为零。用示波器 AC 耦合观察，看到的是正弦波；用万用表交流挡读有效值 $U_2$，用示波器读峰值 $\sqrt2U_2$ 和周期 $20\,\mathrm{ms}$。波形见 §2.2 与图 2(a)：横轴标 $\omega t$ 或 $t$，纵轴标 $u_2$，标出峰值 $\sqrt2U_2$ 和周期 $T=20\,\mathrm{ms}$。

\subhdr{4.2 桥式整流：输出与副边电压是什么关系？输出是直流还是交流？画出波形}

输出平均值 $U_{o1}\approx0.9U_2$（忽略二极管压降）；输出是\textbf{脉动直流}，方向不变但大小仍在变，所以既不是纯交流，也不是稳定直流。波形见 图 2(b)：把负半周翻到正半周，脉动频率 $100\,\mathrm{Hz}$，周期 $10\,\mathrm{ms}$；每只二极管导通半个周期，电流 $I_D=I_o/2$，承受反向电压 $\sqrt2U_2$。

\subhdr{4.3 电容滤波：输出与副边电压是什么关系？输出是直流还是交流？画出波形}

带负载时 $U_{o1}\approx1.2U_2$；空载时 $U_{o1}\approx\sqrt2U_2\approx1.4U_2$。输出是\textbf{带纹波的直流}，直流分量约为 $1.2U_2$，纹波峰峰值 $V_r\approx I_o/(2fC_1)$。波形见 图 2(c)：整体是一条较高的水平线，上面叠着小的锯齿；电容在整流电压高时充电、低时放电，所以谷值被填高、峰值被保留。

\subhdr{4.4 稳压管输出是直流还是交流？它随负载电阻变化吗？}

输出是\textbf{直流}，在稳压管正常工作时近似为 $U_Z=12\,\mathrm{V}$。它并不是“任何负载下都绝对不变”：在允许的负载范围内，$I_Z$ 变化只引起很小的 $U_{o2}$ 变化（由动态电阻 $r_Z$ 决定）；当负载过重、$I_Z$ 降到 $I_{Z\min}$ 以下时，稳压管退出击穿，$U_{o2}$ 会明显下降；当负载过轻、$I_Z$ 超过 $I_{Z\max}$ 时，稳压管可能过热损坏。所以正确的说法是：\textbf{在稳压范围内近似恒压，超出范围就不再稳压}。

%=============================================================
\hdr{§5}{实验流程：按顺序做，不跳步}

\subhdr{5.1 课前准备（进实验室之前）}

\begin{enumerate}
  \item 读手册 p10--p12，画好实验电路图（标出 $T_1$、$D_1$、$C_1$、$R_1$、$R_2$、$D_Z$、$R_w$、$R_3$ 和三个开关），准备好 Table 1、Table 2 的空表。
  \item 按 §2.6 用你预计的 $U_2$ 先算一遍：$0.9U_2$、$1.2U_2$、$U_{o2}$ 两种负载下的预算值；这样现场读数时能立刻判断“对不对”。
  \item 复习 DS2202A 示波器：探头衰减、耦合方式、垂直/水平挡位、触发方式、光标测量；见 §5.8。
  \item 检查实验台：变压器模块、桥堆、电容、电阻、稳压管、导线、万用表、示波器探头是否齐全；电容有无鼓包、电阻有无烧焦。
\end{enumerate}

\subhdr{5.2 接线与自查（断电状态）}

\begin{enumerate}
  \item 关断变压器模块电源，把调压器/电源开关置于关位。
  \item 按手册 p11 的图接线：先接副边到桥堆，再接桥堆输出到三条支路；$C_1$ 注意极性（正极接正母线），$D_Z$ 注意阴极接正母线、阳极接负母线。
  \item 三个开关先全部断开；$R_w$ 先置于中间值，不要一上来就 0 或 100\,k$\Omega$。
  \item 自查：用万用表电阻挡（断电）检查正、负母线之间没有直接短路；检查 $C_1$、$D_Z$ 极性；检查示波器地线夹准备接负母线。
  \item 请老师检查接线，确认后再通电。
\end{enumerate}

\subhdr{5.3 步骤 1：测变压器副边 $U_2$}

\begin{enumerate}
  \item 保持 $K_1$、$K_2$、$K_3$ 全断，接通变压器模块电源。
  \item 万用表置\textbf{交流电压挡}，量副边两端，读 $U_2$；若读数很小或为零，先检查副边接线和保险丝。
  \item 示波器 CH1 接副边（AC 耦合，探头 $\times10$），调垂直挡位使波形占 4--6 格，水平挡位使屏上出现 2--3 个周期；读峰值、周期、频率。
  \item 记录：$U_2$（有效值）、峰值、周期、频率；在 Table 1 的波形栏画出正弦波，标出 $\sqrt2U_2$、$T=20\,\mathrm{ms}$。
\end{enumerate}

\subhdr{5.4 步骤 2：桥式整流（只合 $K_2$）}

\begin{enumerate}
  \item 关电源，合 $K_2$，断 $K_1$、$K_3$；再通电。
  \item 万用表置\textbf{直流电压挡}，量 $R_1$ 两端，读 $U_{o1}$；同时用交流挡量一次，交流分量应明显小于直流分量（说明是脉动直流，不是纯交流）。
  \item 示波器改 DC 耦合，观察 $R_1$ 两端波形：应是连续两个正半波，谷点接近零，脉动周期 $10\,\mathrm{ms}$。
  \item 记录 $U_{o1}$，画波形；标注峰值、谷值、周期 $10\,\mathrm{ms}$。
\end{enumerate}

\subhdr{5.5 步骤 3：电容滤波（合 $K_1$、$K_2$）}

\begin{enumerate}
  \item 关电源，合 $K_1$（此时 $K_2$ 已合，$K_3$ 断）；再通电。
  \item 万用表直流挡量 $R_1$ 两端，读滤波后的 $U_{o1}$；读数应比步骤 2 明显升高（约 $1.2U_2$）。
  \item 示波器 DC 耦合观察：波形应是一条较高的直流线，上面有小的锯齿；若纹波太大，检查 $C_1$ 是否接好、负载是否过重。
  \item 若想看清纹波，把示波器改 AC 耦合、垂直挡位调小（如 $0.2\,\mathrm{V/div}$），读纹波峰峰值 $V_r$；再换回 DC 耦合读直流分量。
  \item 记录 $U_{o1}$、$V_r$，画波形；标注直流分量、纹波峰峰值、周期 $10\,\mathrm{ms}$。
\end{enumerate}

\subhdr{5.6 步骤 4：稳压管（合 $K_1$、$K_3$，断 $K_2$）}

\begin{enumerate}
  \item 关电源，断 $K_2$，合 $K_1$、$K_3$；把 $R_w$ 调到 $0\,\Omega$，此时 $R_w+R_3=3\,\mathrm{k}\Omega$。
  \item 通电，万用表直流挡量 $R_w+R_3$ 两端（即 $U_{o2}$），读第一组 $U_{o2}$；同时量 $R_2$ 两端电压，便于算 $I_R$。
  \item 关电源，把 $R_w$ 调到 $100\,\mathrm{k}\Omega$，此时 $R_w+R_3=103\,\mathrm{k}\Omega$；通电，读第二组 $U_{o2}$。
  \item 示波器 DC 耦合观察 $U_{o2}$：若稳压管工作，应是一条较平的直流线；若在 $3\,\mathrm{k}\Omega$ 时看到明显下凹，说明稳压管已退出击穿。
  \item 记录两组 $U_{o2}$，填 Table 2；用 §2.5 的公式判断每组稳压管是否在工作。
\end{enumerate}

\subhdr{5.7 关机与收尾}

\begin{enumerate}
  \item 先关变压器模块电源，再关示波器、万用表；拆线前用电阻给 $C_1$ 放电。
  \item 把 $R_w$ 调回中间值，开关复位，导线整理好；不要把电容、稳压管随手放在金属台面上。
  \item 请老师在原始数据上签字/批改；未批改的记录无效（手册 p02 的上课须知）。
  \item 若数据异常，先记录现象（哪个开关、哪个挡位、读数多少），再报告老师；不要事后编数据。
\end{enumerate}

\subhdr{5.8 DS2202A 示波器速查}

\begin{center}
\small
\begin{tabular}{@{}p{3.4cm}p{1.5cm}p{3.2cm}p{5.0cm}@{}}
\hline
\thead{测量对象} & \thead{耦合} & \thead{垂直/水平} & \thead{触发/备注} \\
\hline
$U_2$（约 14\,V） & AC & 5\,V/div，5\,ms/div & CH1 边沿、Auto；探头 $\times10$ \\
整流 $U_{o1}$ & DC & 5\,V/div，5\,ms/div & 看谷点是否到零 \\
滤波 $U_{o1}$ & DC & 5\,V/div，5\,ms/div & 纹波细节：AC，0.2\,V/div \\
稳压 $U_{o2}$ & DC & 5\,V/div，5\,ms/div & 看是否平直、是否下凹 \\
\hline
\end{tabular}
\end{center}

\aside{DS2202A 是双通道示波器，两个通道共地。测低压侧时，探头地线夹接负母线；不要把地线夹接到正母线或开关两端，否则会通过示波器电源地造成短路。探头衰减开关要与示波器菜单里的探头倍率一致（都用 $\times10$），否则读数会差 10 倍。}

%=============================================================
\hdr{§6}{数据记录表与预算}

\subhdr{6.1 Table 1：变压器、整流、滤波}

\begin{center}
\small
\begin{tabular}{@{}p{4.4cm}p{3.0cm}p{3.0cm}p{5.4cm}@{}}
\hline
\thead{参数} & \thead{计算值/V} & \thead{测量值/V} & \thead{输出波形（画 2 个周期）} \\
\hline
变压器副边 $U_2$ & 铭牌值 & \hspace{2.4cm} & 正弦，峰值 $\sqrt2U_2$，$T=20\,\mathrm{ms}$ \\
桥式整流 $U_{o1}$ & $0.9U_2$ & & 全波脉动，谷点近零，$T=10\,\mathrm{ms}$ \\
电容滤波 $U_{o1}$ & $1.2U_2$ & & 带纹波直流，$T=10\,\mathrm{ms}$ \\
\hline
\end{tabular}
\end{center}

\subhdr{6.2 Table 2：稳压输出}

\begin{center}
\small
\begin{tabular}{@{}p{4.2cm}p{4.0cm}p{4.0cm}p{3.6cm}@{}}
\hline
\thead{$R_w+R_3$} & \thead{计算值 $U_{o2}$/V} & \thead{测量值 $U_{o2}$/V} & \thead{稳压管状态} \\
\hline
$3\,\mathrm{k}\Omega$ & 先假设 $12\,\mathrm{V}$，算 $I_Z$；若 $I_Z<0$，改 $0.6U_{o1}$ & & 工作/退出 \\
$103\,\mathrm{k}\Omega$ & 先假设 $12\,\mathrm{V}$，算 $I_Z$ & & 工作/退出 \\
\hline
\end{tabular}
\end{center}

\subhdr{6.3 预算示例：$U_2=14.0\,\mathrm{V}$ 时两张表怎么填}

若测得 $U_2=14.0\,\mathrm{V}$，则 Table 1 的计算值为：整流 $0.9\times14.0=12.6\,\mathrm{V}$，滤波 $1.2\times14.0=16.8\,\mathrm{V}$。Table 2 先按 $U_{o1}=16.8\,\mathrm{V}$ 算：
\[
I_R=\frac{16.8-12}{2\,\mathrm{k}\Omega}=2.40\,\mathrm{mA}.
\]
$3\,\mathrm{k}\Omega$ 时，若稳压管工作，$I_L=12/3\,\mathrm{k}=4.00\,\mathrm{mA}$，$I_Z=2.40-4.00=-1.60\,\mathrm{mA}<0$，说明假设不成立，稳压管退出；改按分压
\[
U_{o2}=\frac{3}{2+3}\times16.8=10.08\,\mathrm{V}.
\]
$103\,\mathrm{k}\Omega$ 时，$I_L=12/103\,\mathrm{k}=0.1165\,\mathrm{mA}$，$I_Z=2.40-0.1165=2.28\,\mathrm{mA}>0$，稳压管工作，$U_{o2}=12\,\mathrm{V}$。这两个数与 Table 2 的“计算值”栏对应；测量值以现场读数为准，误差通常来自 $U_{o1}$ 的实际值、电阻容差和万用表精度。

\subhdr{6.4 记录规范}

\begin{itemize}
  \item 每个数后面写单位（V、mA、$\mathrm{k}\Omega$、$\mu\mathrm{F}$）；没有单位的数在报告里等于错。
  \item 波形图要有坐标轴、纵轴量名与单位、横轴量名与单位、周期/峰值标注；只画一条曲线不标轴，最多拿一半分。
  \item 计算值、测量值、相对误差分栏写；相对误差用 $\delta=|X_{\mathrm{meas}}-X_{\mathrm{calc}}|/|X_{\mathrm{calc}}|\times100\%$。
  \item 记录原始读数，不要只写“约 12\,V”；至少保留三位有效数字（如 12.6\,V、16.8\,V、2.40\,mA）。
\end{itemize}

%=============================================================
\hdr{§7}{数据分析：五个必须会算的量}

\subhdr{7.1 输出/输入比：整流约 0.9，滤波约 1.2}

用 Table 1 的同一组数据，分别算
\[
r_{\text{整流}}=\frac{U_{o1,\text{整流}}}{U_2},\qquad
r_{\text{滤波}}=\frac{U_{o1,\text{滤波}}}{U_2}.
\]
理论值分别是 0.9 与 1.2。算相对误差：
\[
\delta_{\text{整流}}=\frac{|r_{\text{整流}}-0.9|}{0.9}\times100\%,\qquad
\delta_{\text{滤波}}=\frac{|r_{\text{滤波}}-1.2|}{1.2}\times100\%.
\]
若 $U_2=14.0\,\mathrm{V}$、整流测得 $12.5\,\mathrm{V}$、滤波测得 $16.6\,\mathrm{V}$，则 $r_{\text{整流}}=0.893$，$\delta=0.8\%$；$r_{\text{滤波}}=1.186$，$\delta=1.2\%$。这两个误差都在正常范围内。若滤波比明显低于 1.2，先检查 $C_1$ 是否接好、负载是否过重、万用表是否用了直流挡；若明显高于 1.4，检查是不是空载或 $C_1$ 容量偏大。

\subhdr{7.2 若 $C_1$ 短路，会发生什么？}

$C_1$ 短路后，桥堆的输出正、负母线被直接短接，$U_{o1}=0$，滤波作用完全消失。更危险的是：正半周 $D_1$、$D_4$ 导通，负半周 $D_2$、$D_3$ 导通，电流不经过负载，只经过“两只二极管 + 短路线”，回路里只剩变压器绕组电阻和二极管正向电阻，电流会很大。后果可能是：整流二极管过热损坏、变压器绕组过热、保险丝熔断，甚至 $C_1$ 本身因短路电流发热。所以实验中不要用导线或螺丝刀跨接 $C_1$；若要放电，用电阻，不要直接短路。

\subhdr{7.3 哪个电阻是限流电阻？若它等于零会怎样？}

限流电阻是 $R_2=2\,\mathrm{k}\Omega$，它串在滤波输出与稳压管之间。正常时它承担 $U_{o1}-U_Z$ 的电压差，把 $I_Z$ 限制在安全范围：
\[
I_R=\frac{U_{o1}-U_Z}{R_2}.
\]
若 $R_2=0$，稳压管直接并在滤波输出两端，$I_Z$ 只受变压器内阻、二极管正向电阻和导线电阻限制；只要 $U_{o1}>U_Z$，稳压管就会流过很大的电流，超过 $I_{Z\max}$ 而损坏，严重时还会拖垮整流桥或变压器。所以 $R_2$ 不能短接，也不能随意换小；要改稳压值或负载范围，应重新算 $R_2$ 的阻值和功率。

\subhdr{7.4 稳压管在 $3\,\mathrm{k}\Omega$ 时到底工不工作？}

按 §2.5 的三步法算，不要先看测量值“像不像 12\,V”。以 $U_{o1}=16.8\,\mathrm{V}$ 为例：
\[
I_R=\frac{16.8-12}{2\,\mathrm{k}\Omega}=2.40\,\mathrm{mA},\quad
I_L=\frac{12}{3\,\mathrm{k}\Omega}=4.00\,\mathrm{mA},
\]
\[
I_Z=I_R-I_L=2.40-4.00=-1.60\,\mathrm{mA}<0.
\]
$I_Z$ 为负，说明“稳压管工作”的假设不成立：负载把 $R_2$ 送来的电流全部拿走还不够，稳压管不可能还有电流。此时它退出击穿，输出由 $R_2$ 与 $R_L$ 分压：
\[
U_{o2}=\frac{3}{2+3}\times16.8=10.08\,\mathrm{V}.
\]
所以 $3\,\mathrm{k}\Omega$ 这一档的答案是“\textbf{不工作}”，输出约 $10\,\mathrm{V}$，不是 $12\,\mathrm{V}$。到了 $103\,\mathrm{k}\Omega$，$I_L=0.1165\,\mathrm{mA}$，$I_Z=2.28\,\mathrm{mA}>0$，稳压管工作，$U_{o2}\approx12\,\mathrm{V}$。把判据写成一般式更清楚：
\[
I_Z<0\ \Longleftrightarrow\ U_{o1}<U_Z\left(1+\frac{R_2}{R_L}\right).
\]
对 $R_L=3\,\mathrm{k}\Omega$，临界值是 $U_Z(1+2/3)=20\,\mathrm{V}$；只要测得的 $U_{o1}$ 低于 $20\,\mathrm{V}$，这一档的稳压管必然退出。在本例 $U_{o1}=16.8\,\mathrm{V}$ 下，这不是“碰巧不工作”，而是由 $U_{o1}$ 与 $R_2/R_L$ 共同决定的。这也说明：稳压管不是“接上就稳压”，它的工作范围由 $R_2$、$U_{o1}$ 和负载共同决定。

\subhdr{7.5 副边波形是不是标准正弦？}

不是“标准”正弦，只能说是\textbf{近似正弦}。原因有三类：（1）电网电压本身可能有畸变；（2）变压器有励磁电流、绕组电阻和漏抗，带负载时波形会平顶或畸变；（3）整流桥和滤波电容是非线性负载，电流只在部分时间流通，会反过来影响副边电压。示波器上我们主要读它的周期、频率和峰值，不要因为“看起来不是完美正弦”就以为变压器坏了。判断方法：周期是否 $20\,\mathrm{ms}$、频率是否 $50\,\mathrm{Hz}$、峰值是否约 $\sqrt2U_2$；只要这三个量对，变压器就是正常的。

\subhdr{7.6 误差从哪里来？}

\begin{center}
\small
\begin{tabular}{@{}p{3.8cm}p{6.2cm}p{5.4cm}@{}}
\hline
\thead{误差来源} & \thead{表现} & \thead{减小办法} \\
\hline
万用表精度/挡位 & 交直流挡读数不同，末位跳动 & 选合适量程；读稳定值；记录挡位 \\
电阻容差 & $R_1$、$R_2$、$R_3$ 偏离标称值 & 用实测电阻值代入计算 \\
二极管正向压降 & 整流输出比 $0.9U_2$ 略低 & 在结论里说明“忽略二极管压降” \\
电容容差/漏电 & 滤波输出偏离 $1.2U_2$ & 用实测 $C_1$ 和 $U_{o1}$ 分析 \\
纹波与负载效应 & 直流挡读数含纹波平均 & 同时读纹波，说明取的是平均值 \\
示波器读数 & 峰峰值、周期读不准 & 用光标/自动测量；探头倍率一致 \\
\hline
\end{tabular}
\end{center}

%=============================================================
\hdr{§8}{Discussion 五题参考答法（手册 p13）}

\subhdr{8.1 第 1 题：输出/输入比与理论值比较}

先写公式，再代入 Table 1 的数：
\[
r_{\text{整流}}=\frac{U_{o1,\text{整流}}}{U_2},\qquad
r_{\text{滤波}}=\frac{U_{o1,\text{滤波}}}{U_2}.
\]
理论值：整流 $0.9$，滤波 $1.2$（带负载）。以 $U_2=14.0\,\mathrm{V}$、整流 $12.5\,\mathrm{V}$、滤波 $16.6\,\mathrm{V}$ 为例：
\[
r_{\text{整流}}=12.5/14.0=0.893,\quad \delta=0.8\%;
\]
\[
r_{\text{滤波}}=16.6/14.0=1.186,\quad \delta=1.2\%.
\]
结论：两个比值都与理论值吻合，说明整流、滤波电路工作正常；滤波比略低于 1.2 是负载、电容容差和二极管压降造成的。若你的数据偏离较大，按 §7.1 的顺序排查。

\subhdr{8.2 第 2 题：$C_1$ 短路会怎样}

$C_1$ 短路后，正负母线直接短接，$U_{o1}=0$，滤波消失；整流二极管在每个半周轮流导通，电流不经过负载，回路电阻很小，电流很大，可能烧坏二极管、变压器绕组或保险丝。回答时先写“输出为零、滤波失效”，再写“桥堆经短路支路导通、电流过大、可能损坏元件”，最后写“实验中禁止直接短接电容，放电要用电阻”。

\subhdr{8.3 第 3 题：哪个是限流电阻？它为零会怎样}

限流电阻是 $R_2=2\,\mathrm{k}\Omega$。它串在 $U_{o1}$ 与稳压管之间，正常时 $I_R=(U_{o1}-U_Z)/R_2$。若 $R_2=0$，$I_Z$ 只由电源内阻和二极管正向电阻限制，会超过 $I_{Z\max}$，稳压管过热损坏，甚至拖垮整流桥。所以 $R_2$ 既不能短接，也不能随意减小；它的功率要按最坏情况 $P_{R2}=I_R^2R_2$ 选。

\subhdr{8.4 第 4 题：$R_w+R_3=3\,\mathrm{k}\Omega$ 时稳压管是否工作}

按 §7.4 的三步法：先假设 $U_{o2}=12\,\mathrm{V}$，算 $I_R=2.40\,\mathrm{mA}$、$I_L=4.00\,\mathrm{mA}$，得 $I_Z=-1.60\,\mathrm{mA}<0$，假设不成立；稳压管退出，按分压 $U_{o2}=0.6U_{o1}\approx10.1\,\mathrm{V}$。结论：在 $U_{o1}=16.8\,\mathrm{V}$ 的例子中，$3\,\mathrm{k}\Omega$ 时稳压管\textbf{不工作}，输出约 $10\,\mathrm{V}$；$103\,\mathrm{k}\Omega$ 时 $I_Z=2.28\,\mathrm{mA}>0$，稳压管工作，输出 $12\,\mathrm{V}$。这道题的关键不是背“3\,k 不工作”，而是会算 $I_Z$ 并会用 $I_Z$ 的正负判断；若你的 $U_{o1}$ 高于 $20\,\mathrm{V}$，$3\,\mathrm{k}\Omega$ 这一档也可能稳压，所以报告里要写清你代入的 $U_{o1}$。

\subhdr{8.5 第 5 题：副边波形是不是标准正弦}

不是标准正弦，是近似正弦。原因：电网畸变、变压器励磁电流与绕组阻抗、整流/滤波的非线性负载。判断变压器是否正常，看周期 $20\,\mathrm{ms}$、频率 $50\,\mathrm{Hz}$、峰值约 $\sqrt2U_2$ 是否对，而不是看波形“完美不完美”。若波形严重削顶、周期不对或峰值明显偏低，再检查变压器、接线和示波器探头。

%=============================================================
\hdr{§9}{常见错误与评分点}

\begin{center}
\small
\begin{tabular}{@{}p{4.2cm}p{5.6cm}p{5.2cm}@{}}
\hline
\thead{常见错误} & \thead{后果} & \thead{正确做法} \\
\hline
$K_1$/$K_2$/$K_3$ 投错 & 读数为 0、波形不对、误判故障 & 按 §0.3 真值表逐项核对 \\
交直流挡用错 & $U_2$ 读成 0，$U_{o1}$ 读数偏大/偏小 & $U_2$ 用 AC，$U_{o1}$/$U_{o2}$ 用 DC \\
示波器地线夹错位置 & 短路、跳闸、波形异常 & 地线夹接负母线 \\
$C_1$ 极性接反 & 电容发热、鼓包、爆裂 & 正极接正母线，通电前核对 \\
带电改接 & 触电、损坏元件 & 先关电源，放电后再改 \\
波形不标轴/周期/峰值 & 波形题丢分 & 标 $u$、$t$、$T$、$U_m$ \\
计算不写公式 & 计算题丢分 & 公式—代入—结果—结论四步 \\
用标称值代替实测值 & 误差分析对不上 & 用实测 $U_2$、$R$、$C$ \\
\hline
\end{tabular}
\end{center}

评分结构（手册 p01）：Preparation 5 分、Operation 5 分、Discussion 10+5 分。预习报告、接线规范、数据记录、波形图、计算过程都算分；其中“写出公式并代入数值”是总要求，适用于所有计算题。

%=============================================================
\hdr{§10}{公式卡与 60 秒自检}

\subhdr{10.1 公式卡}

\[
U_{o1}=0.9U_2\ \text{（桥式整流）},\qquad
U_{o1}=1.2U_2\ \text{（电容滤波，带负载）},\qquad
U_{o1}=\sqrt2U_2\ \text{（空载）}.
\]
\[
I_o=\frac{U_{o1}}{R_L},\qquad
I_D=\frac{I_o}{2},\qquad
U_{Rm}=\sqrt2U_2.
\]
\[
V_r\approx\frac{I_o}{2fC_1},\qquad
I_R=\frac{U_{o1}-U_Z}{R_2},\qquad
I_L=\frac{U_Z}{R_L},\qquad
I_Z=I_R-I_L.
\]
\[
\text{稳压管工作： }I_{Z\min}\le I_Z\le I_{Z\max};\qquad
\text{退出： }I_Z<0\ \Rightarrow\ U_{o2}=\frac{R_L}{R_2+R_L}U_{o1}.
\]

\subhdr{10.2 交报告前 60 秒自检}

\begin{enumerate}
  \item Table 1、Table 2 的每个数都有单位；计算值、测量值、相对误差分栏。
  \item 三个波形都画了坐标轴、周期、峰值/谷值；整流波形谷点接近零，滤波波形有纹波。
  \item 五道 Discussion 每题都有“公式—代入—结果—结论”，没有只写答案。
  \item $3\,\mathrm{k}\Omega$ 那一档写清了“稳压管不工作”的理由（$I_Z<0$），不是只写 $10\,\mathrm{V}$。
  \item 示波器探头倍率、万用表挡位、开关状态都写进了报告或原始记录。
  \item 原始数据有老师签字/批改；没有编数据、抄数据。
\end{enumerate}

%=============================================================
\hdr{§11}{三句话总结}

\begin{enumerate}
  \item 桥式整流把交流的负半周翻正，输出平均值约 $0.9U_2$；电容滤波把谷值填高，带负载时约 $1.2U_2$，空载时接近 $1.4U_2$。
  \item 稳压管不是“接上就稳压”：先用 $I_R=(U_{o1}-U_Z)/R_2$、$I_L=U_Z/R_L$ 算 $I_Z$，只有 $I_Z$ 在允许范围内才稳压；$3\,\mathrm{k}\Omega$ 时 $I_Z<0$，它退出，输出约 $0.6U_{o1}$。
  \item 实验的关键不是“读到 12\,V”，而是会用波形、比值和 $I_Z$ 判断每一级是否正常工作；所有计算都要写公式、代数值、给结论。
\end{enumerate}

%=============================================================
\hdr{§12}{页码对照与缺口清单}

\subhdr{12.1 本讲义与手册/课件的对应}

\begin{center}
\small
\begin{tabular}{@{}p{2.2cm}p{6.0cm}p{6.4cm}@{}}
\hline
\thead{手册页} & \thead{内容} & \thead{课件/教材依据} \\
\hline
p10 & 预习 4 问 & C3-2 pp.13--25；C3-3 pp.8--11、p15--24 \\
p11 & 电路与整流/滤波步骤 & C3-2 p23（桥式）、pp.24--25（电容滤波） \\
p12 & Table 1/Table 2、稳压步骤 & C3-3 pp.9--11（$0.9/1.2/\sqrt2$）、pp.19--24（稳压管） \\
p13 & Discussion 5 问 & C3-3 p11、pp.19--24；本讲义 §7--§8 \\
\hline
\end{tabular}
\end{center}

\subhdr{12.2 手册没有明说、本讲义补上的五处}

\gapitem{开关状态矛盾}{手册 p11 第三步写“disconnect K1 and K2, connect K3”却又要求量 $5.1\,\mathrm{k}\Omega$ 上的 $U_{o1}$；按电路，这一步应是 $K_1$、$K_2$ 合，$K_3$ 断。§0.3 给出真值表并说明理由。}
\gapitem{滤波公式的条件}{手册 p10 只问“under load”的输出关系，没给 $1.2U_2$ 的适用条件；§2.4 补了带负载 $1.2U_2$、空载 $1.4U_2$、纹波估算 $V_r\approx I_o/(2fC_1)$。}
\gapitem{稳压管退出后的算法}{手册 p12 的 Table 2 只留“calculated value”，没说 $I_Z<0$ 时怎么算；§2.5 与 §7.4 补了三步判据和分压式 $U_{o2}=R_LU_{o1}/(R_2+R_L)$。}
\gapitem{“是否工作”的判据}{手册 p13 第 4 问只问“calculate whether the Zener is in operation”，没给判据；§7.4 用 $I_Z=I_R-I_L$ 的正负和范围来判断，而不是看输出像不像 12\,V。}
\gapitem{波形参数}{手册 p10 要求“draw the output waveform”，但没规定标什么；§5.3--§5.6 与 §6.4 明确要求标坐标轴、周期、峰值/谷值、直流分量与纹波。}

%=============================================================
\hdr{§13}{一页纸流程卡（可单独打印）}

\begin{center}
\small
\begin{tabular}{@{}p{1.0cm}p{3.0cm}p{4.2cm}p{7.0cm}@{}}
\hline
\thead{步} & \thead{开关状态} & \thead{测什么} & \thead{预期/要点} \\
\hline
1 & 全断 & 副边 $U_2$ & 正弦；$U_2$ 用 AC 挡；峰值 $\sqrt2U_2$，$T=20\,\mathrm{ms}$ \\
2 & $K_2$ 合 & $R_1$ 两端 $U_{o1}$ & 全波脉动，$U_{o1}\approx0.9U_2$，$T=10\,\mathrm{ms}$ \\
3 & $K_1$、$K_2$ 合 & $R_1$ 两端 $U_{o1}$ & 带纹波直流，$U_{o1}\approx1.2U_2$，$T=10\,\mathrm{ms}$ \\
4 & $K_1$、$K_3$ 合，$K_2$ 断，$R_w=0$ & $R_w+R_3$ 两端 $U_{o2}$ & $R_L=3\,\mathrm{k}\Omega$；算 $I_Z$，若 $I_Z<0$ 则稳压管退出，$U_{o2}\approx0.6U_{o1}$ \\
5 & 同上，$R_w=100\,\mathrm{k}\Omega$ & $U_{o2}$ & $R_L=103\,\mathrm{k}\Omega$；$I_Z>0$，$U_{o2}\approx12\,\mathrm{V}$ \\
\hline
\end{tabular}
\end{center}

\warnbox{\textbf{三条最容易丢分的点。} （1）$K_1$--$K_3$ 投错：按上表逐项核对，不要凭“表上有数”判断；（2）$3\,\mathrm{k}\Omega$ 那一档只写“约 10\,V”不写理由：必须写出 $I_R$、$I_L$、$I_Z<0$ 和分压式；（3）波形图不标轴、不标周期/峰值：每个波形至少标出纵轴量、横轴量、周期和一个峰值/谷值。}

\begin{center}
\small
\begin{tabular}{@{}p{8.4cm}p{8.0cm}@{}}
\hline
\thead{核心公式} & \thead{判据/结论} \\
\hline
$U_{o1}=0.9U_2$（桥式整流） & 全波整流平均值；二极管压降忽略时成立 \\
$U_{o1}=1.2U_2$（带负载滤波） & 空载时 $\sqrt2U_2\approx1.4U_2$ \\
$V_r\approx I_o/(2fC_1)$ & $I_o=U_{o1}/R_L$；纹波随 $C_1$、$R_L$ 增大而减小 \\
$I_R=(U_{o1}-U_Z)/R_2$ & $R_2$ 是限流电阻，不能短接 \\
$I_L=U_Z/R_L$，$I_Z=I_R-I_L$ & $I_{Z\min}\le I_Z\le I_{Z\max}$ 才稳压 \\
$U_{o2}=R_LU_{o1}/(R_2+R_L)$ & $I_Z<0$ 时稳压管退出，按分压算 \\
\hline
\end{tabular}
\end{center}
\end{document}
